\clearpage\pagestyle{otMain}\setlength{\otBlocEspacement}{2.5mm}\fontsize{11}{15}\selectfont\corrigesChapitre{7}{Carrés et cubes}
\begin{blocLivre}{cor-7-1}\begin{corrige}{1}{Lire un carré}\demarche{Déplier la notation.}
$9^2=9\times9=\res{81}$. Il y a deux facteurs 9.
\end{corrige}\end{blocLivre}
\begin{blocLivre}{cor-7-2}\begin{corrige}{2}{Lire un cube}\demarche{Compter trois facteurs.}
$4^3=4\times4\times4=16\times4=\res{64}$.
\end{corrige}\end{blocLivre}
\begin{blocLivre}{cor-7-3}\begin{corrige}{3}{Carrés repères}\demarche{Utiliser les carrés à connaître.}
Ils valent $\res{0,\ 49,\ 121,\ 144}$. Ces valeurs se retrouvent dans les produits des nombres par eux-mêmes.
\end{corrige}\end{blocLivre}
\begin{blocLivre}{cor-7-4}\begin{corrige}{4}{Retrouver la base}\demarche{Reconnaître les produits répétés.}
$6\times6=36$ et $10\times10\times10=1000$. Les bases sont $\res6$ et $\res{10}$.
\end{corrige}\end{blocLivre}
\begin{blocLivre}{cor-7-5}\begin{corrige}{5}{Puissance ou produit}\demarche{Distinguer exposant et facteur.}
$5^3=\res{125}$ tandis que $5\times3=\res{15}$. Les écritures ne décrivent pas le même calcul.
\end{corrige}\end{blocLivre}
\begin{blocLivre}{cor-7-6}\begin{corrige}{6}{Priorités avec un carré}\demarche{Calculer la puissance avant le produit.}
Le carré vaut 16, puis le produit vaut 48. Le total est $\res{54}$.
\end{corrige}\end{blocLivre}
\begin{blocLivre}{cor-7-7}\begin{corrige}{7}{Parenthèses et carré}\demarche{Identifier ce qui est élevé au carré.}
Le premier vaut $5^2=\res{25}$. Le second vaut $64-9=\res{55}$. Les parenthèses définissent la base.
\end{corrige}\end{blocLivre}
\begin{blocLivre}{cor-7-8}\begin{corrige}{8}{Carré décimal}\demarche{Multiplier les facteurs identiques.}
$1{,}5\times1{,}5=\res{2{,}25}$. $0{,}2\times0{,}2\times0{,}2=\res{0{,}008}$. Une puissance peut diminuer un nombre compris entre 0 et 1.
\end{corrige}\end{blocLivre}
\begin{blocLivre}{cor-7-9}\begin{corrige}{9}{Une substitution}\demarche{Remplacer la lettre avec soin.}
$A=3\times25+2=\res{77}$ et $B=8-2=\res6$. Le coefficient 3 ne fait pas partie de la base du carré.
\end{corrige}\end{blocLivre}
\begin{blocLivre}{cor-7-10}\begin{corrige}{10}{Une unité de volume}\demarche{Convertir les trois dimensions.}
Un cube de 1 dm de côté a 10 cm de côté : $10^3=1000$ cm$^3$. Donc 2 dm$^3=\res{2000\ \mathrm{cm}^3}$.
\end{corrige}\end{blocLivre}
\begin{blocLivre}{cor-7-11}\begin{corrige}{11}{Un carré à carreler}\demarche{Comparer des aires dans la même unité.}
L’aire vaut $3{,}2^2=\res{10{,}24\ \mathrm{m}^2}$. Quatre paquets couvrent 10 m$^2$, ce qui ne suffit pas : il faut cinq paquets au minimum, en négligeant les pertes.
\label{fig-cor-7-11}\figureGeo[corrige]{puiss-11}
\end{corrige}\end{blocLivre}
\begin{blocLivre}{cor-7-12}\begin{corrige}{12}{Un cube de cubes}\demarche{Distinguer couche et solide entier.}
Une couche contient $6^2=\res{36}$ cubes. Il y a six couches, donc $6^3=\res{216}$ cubes.
\label{fig-cor-7-12}\figureGeo[corrige]{puiss-12}
\end{corrige}\end{blocLivre}
\begin{blocLivre}{cor-7-13}\begin{corrige}{13}{Deux cubes comparés}\demarche{Comparer les produits des trois dimensions.}
Les volumes valent $2^3=8$ et $4^3=64$ cm$^3$. Le second vaut $\res8$ fois le premier. Doubler chacune des trois dimensions donne $2\times2\times2=8$.
\label{fig-cor-7-13}\figureGeo[corrige]{puiss-13}
\end{corrige}\end{blocLivre}
\begin{blocLivre}{cor-7-14}\begin{corrige}{14}{Réfuter une formule}\demarche{Tester des valeurs simples.}
Pour $a=1$ et $b=2$, $(1+2)^2=9$, tandis que $1^2+2^2=5$. Comme $\res{9\ne5}$, l’affirmation générale est fausse. Un seul contre-exemple suffit.
\end{corrige}\end{blocLivre}
\begin{blocLivre}{cor-7-15}\begin{corrige}{15}{Une bordure carrée}\demarche{Soustraire l’aire intérieure à l’aire extérieure.}
Le côté extérieur vaut $8+1+1=10$ cm. L’aire cherchée est $10^2-8^2=100-64=\res{36\ \mathrm{cm}^2}$. Ajouter seulement 1 au côté oublierait l’autre bord.
\label{fig-cor-7-15}\figureGeo[corrige]{puiss-15}
\end{corrige}\end{blocLivre}
